% ============================================================
\chapter{Applications and Ethics}
\label{ch:applications}
% ============================================================

\begin{intuition}
Deep learning is not an end in itself: it is a tool in the service of science, medicine, industry, and society. This chapter explores the most notable applications and the fundamental ethical questions that every practitioner must consider.
\end{intuition}

% ------------------------------------------------------------
\section{Computer Vision}
\label{sec:app_vision}
\index{computer vision}

\subsection{Object Detection}
\index{object detection}
\index{YOLO}

Object detection simultaneously localizes and classifies multiple objects in an image. The YOLO (\emph{You Only Look Once}) approach divides the image into a grid and directly predicts bounding boxes and class probabilities.

For each grid cell $(i,j)$, the model predicts:
\[
\hat{y}_{ij} = (p_{\mathrm{obj}}, b_x, b_y, b_w, b_h, c_1, \ldots, c_K)
\]
where $p_{\mathrm{obj}}$ is the probability that an object is present, $(b_x, b_y, b_w, b_h)$ define the bounding box, and $c_k$ are the class probabilities.

\subsection{Semantic Segmentation}
\index{segmentation}

Segmentation assigns a class to each pixel. The U-Net\index{U-Net} architecture (Ronneberger et al., 2015) uses an encoder-decoder with skip connections:

\begin{center}
\begin{tikzpicture}[
    block/.style={rectangle, draw=NavyBlue, thick, minimum width=1.2cm, minimum height=0.6cm, fill=blue!10, font=\small},
    skip/.style={->, dashed, gray, thick},
    >=Stealth
]
\node[block] (e1) at (0,0) {Conv 64};
\node[block] (e2) at (1.5,-1) {Conv 128};
\node[block] (e3) at (3,-2) {Conv 256};
\node[block] (bn) at (4.5,-3) {Conv 512};
\node[block] (d3) at (6,-2) {Conv 256};
\node[block] (d2) at (7.5,-1) {Conv 128};
\node[block] (d1) at (9,0) {Conv 64};

\draw[->, thick] (e1) -- (e2);
\draw[->, thick] (e2) -- (e3);
\draw[->, thick] (e3) -- (bn);
\draw[->, thick] (bn) -- (d3);
\draw[->, thick] (d3) -- (d2);
\draw[->, thick] (d2) -- (d1);

\draw[skip] (e1) to[bend left=30] (d1);
\draw[skip] (e2) to[bend left=25] (d2);
\draw[skip] (e3) to[bend left=20] (d3);
\end{tikzpicture}
\end{center}

\subsection{Image Generation}
\index{image generation}

Diffusion models (Chapter~\ref{ch:diffusion}) have revolutionized image generation. Stable Diffusion operates in the latent space of a VAE (Chapter~\ref{ch:vae}), considerably reducing the computational cost.

% ------------------------------------------------------------
\section{Natural Language Processing}
\label{sec:app_nlp}
\index{natural language processing}

\subsection{Language Modeling}
\index{language model}

Large language models (LLMs)\index{LLM} generate text in an auto-regressive manner:
\[
p(\bm{x}) = \prod_{t=1}^T p(x_t \mid x_1, \ldots, x_{t-1}; \theta)
\]

\begin{sota}
\begin{itemize}
\item \textbf{GPT-4} (OpenAI, 2023): multimodal model with advanced reasoning capabilities.
\item \textbf{Claude} (Anthropic, 2024--2025): emphasis on safety and alignment.
\item \textbf{LLaMA 3} (Meta, 2024): high-performing open model.
\item \textbf{Gemini} (Google, 2024): natively multimodal.
\end{itemize}
\end{sota}

\subsection{Machine Translation}

The Transformer\index{Transformer} architecture (Chapter~\ref{ch:transformers}) has enabled major advances in translation, achieving near-human performance on many language pairs.

\subsection{Automatic Summarization and Question Answering}

Pre-trained models such as BERT (bidirectional) and GPT (auto-regressive) are fine-tuned for specific comprehension tasks.

% ------------------------------------------------------------
\section{Scientific Applications}
\label{sec:app_science}
\index{scientific applications}

\subsection{Protein Structure Prediction}
\index{AlphaFold}

AlphaFold 2 (Jumper et al., 2021) solved the 50-year-old protein folding problem. The architecture combines:
\begin{itemize}
\item A Transformer with evolutionary attention over multiple sequence alignments (MSA).
\item A structure module that predicts 3D coordinates.
\item Iterative refinement of predictions.
\end{itemize}

\begin{remark}
AlphaFold 3 (2024) extends prediction to protein-ligand, protein-DNA, and protein-RNA complexes, unifying the modeling of molecular interactions.
\end{remark}

\subsection{Weather Forecasting}
\index{GraphCast}

GraphCast (Lam et al., 2023) uses graph neural networks for global weather forecasting, surpassing traditional numerical models while being 1000$\times$ faster.

\subsection{Drug Discovery}

Generative models explore chemical space to propose new candidate molecules:
\begin{itemize}
\item Molecule generation via VAE or diffusion models.
\item Molecular affinity prediction via GNN (Graph Neural Networks).
\item Multi-objective optimization (efficacy, toxicity, synthesizability).
\end{itemize}

% ------------------------------------------------------------
\section{Timeline of Major Advances}

\begin{center}
\begin{tikzpicture}[
    event/.style={rectangle, draw=NavyBlue, thick, rounded corners, fill=blue!10, text width=3.5cm, font=\small, align=center},
    >=Stealth
]
\draw[thick, ->] (0,0) -- (14,0);

\foreach \x/\yr in {0.5/2012, 2.5/2014, 4.5/2015, 6.5/2017, 8.5/2020, 10.5/2022, 12.5/2024} {
    \draw (\x, -0.2) -- (\x, 0.2);
    \node[below, font=\small] at (\x, -0.3) {\yr};
}

\node[event, above] at (0.5, 0.5) {AlexNet\\ImageNet};
\node[event, below] at (2.5, -0.8) {GAN\\VAE};
\node[event, above] at (4.5, 0.5) {ResNet\\BatchNorm};
\node[event, below] at (6.5, -0.8) {Transformer\\Attention};
\node[event, above] at (8.5, 0.5) {GPT-3\\Diffusion};
\node[event, below] at (10.5, -0.8) {ChatGPT\\Stable Diff.};
\node[event, above] at (12.5, 0.5) {GPT-4\\Claude\\Sora};
\end{tikzpicture}
\end{center}

% ------------------------------------------------------------
\section{Interpretability}
\label{sec:interpretability}
\index{interpretability}

\subsection{Saliency Maps}
\index{saliency map}

Saliency maps identify the input regions most influential for the prediction:
\[
S_i = \left|\frac{\partial f_\theta(\bm{x})}{\partial x_i}\right|
\]

\subsection{Grad-CAM}
\index{Grad-CAM}

Grad-CAM (Selvaraju et al., 2017) uses the gradients from the target convolutional layer to produce a heatmap:
\[
\alpha_k^c = \frac{1}{Z}\sum_i \sum_j \frac{\partial y^c}{\partial A^k_{ij}}, \qquad L_{\mathrm{Grad\text{-}CAM}}^c = \relu\!\left(\sum_k \alpha_k^c A^k\right)
\]

\begin{codeblock}[title=\textbf{Grad-CAM in PyTorch}]
\begin{minted}{python}
import torch
import torch.nn.functional as F

class GradCAM:
    def __init__(self, model, target_layer):
        self.model = model
        self.gradients = None
        self.activations = None
        target_layer.register_forward_hook(
            lambda m, i, o: setattr(self, 'activations', o.detach()))
        target_layer.register_full_backward_hook(
            lambda m, gi, go: setattr(self, 'gradients', go[0].detach()))

    def __call__(self, x, class_idx=None):
        output = self.model(x)
        if class_idx is None:
            class_idx = output.argmax(1)
        self.model.zero_grad()
        one_hot = torch.zeros_like(output)
        one_hot[0, class_idx] = 1
        output.backward(gradient=one_hot)

        weights = self.gradients.mean(dim=[2, 3], keepdim=True)
        cam = F.relu((weights * self.activations).sum(1, keepdim=True))
        cam = F.interpolate(cam, size=x.shape[2:],
                           mode='bilinear', align_corners=False)
        cam = cam / (cam.max() + 1e-8)
        return cam.squeeze().detach().numpy()
\end{minted}
\end{codeblock}

\subsection{SHAP and Shapley Values}
\index{SHAP}

Shapley values, from cooperative game theory, assign to each feature a fair contribution to the prediction:
\[
\phi_i = \sum_{S \subseteq \{1,\ldots,d\} \setminus \{i\}} \frac{|S|!(d-|S|-1)!}{d!}\left[f(S \cup \{i\}) - f(S)\right]
\]

% ------------------------------------------------------------
\section{Bias and Fairness}
\label{sec:fairness}
\index{bias}
\index{fairness}

\subsection{Sources of Bias}

\begin{warning}
Deep learning models can amplify biases present in the training data:
\begin{itemize}
\item \textbf{Representation bias}: certain groups underrepresented in the data.
\item \textbf{Labeling bias}: annotations reflecting human prejudices.
\item \textbf{Aggregation bias}: a single model for heterogeneous populations.
\end{itemize}
\end{warning}

\subsection{Fairness Metrics}

\begin{definition}[Demographic Parity]
A classifier satisfies demographic parity if:
\[
P(\hat{Y} = 1 \mid A = a) = P(\hat{Y} = 1 \mid A = b) \quad \forall a, b
\]
where $A$ is the sensitive attribute.
\end{definition}

\begin{definition}[Equalized Odds]
A classifier satisfies equalized odds if:
\[
P(\hat{Y} = 1 \mid Y = 1, A = a) = P(\hat{Y} = 1 \mid Y = 1, A = b) \quad \forall a, b
\]
\end{definition}

% ------------------------------------------------------------
\section{Privacy}
\label{sec:privacy}
\index{privacy}
\index{differential privacy}

\subsection{Differential Privacy}

\begin{definition}[$(\varepsilon,\delta)$-Differential Privacy]
A randomized mechanism $\mathcal{M}$ satisfies $(\varepsilon,\delta)$-differential privacy if, for all adjacent datasets $D, D'$ (differing by a single example) and every measurable set $S$:
\[
P(\mathcal{M}(D) \in S) \leq e^\varepsilon \cdot P(\mathcal{M}(D') \in S) + \delta
\]
\end{definition}

DP-SGD (Abadi et al., 2016) applies differential privacy to SGD by clipping individual gradients and adding Gaussian noise:
\[
\tilde{\bm{g}}_t = \frac{1}{B}\left(\sum_{i \in \mathcal{B}} \mathrm{clip}(\bm{g}_i, C) + \mathcal{N}(0, \sigma^2 C^2 I)\right)
\]

\subsection{Federated Learning}
\index{federated learning}

Federated learning trains a global model without centralizing the data. Each client $k$ computes a local update, and the server aggregates:
\[
\theta_{t+1} = \sum_{k=1}^K \frac{n_k}{n} \theta_{t+1}^{(k)}
\]

% ------------------------------------------------------------
\section{Environmental Impact}
\label{sec:environment}
\index{environmental impact}

Training large models consumes considerable resources:

\begin{center}
\begin{tabular}{lrr}
\toprule
\textbf{Model} & \textbf{Parameters} & \textbf{Estimated CO$_2$ (tons)} \\
\midrule
BERT-Large & 340M & $\sim$0.6 \\
GPT-3 & 175B & $\sim$500 \\
LLaMA 2 70B & 70B & $\sim$290 \\
GPT-4 (estimated) & $>$1T & $\sim$5\,000+ \\
\bottomrule
\end{tabular}
\end{center}

\begin{bestpractice}
\begin{itemize}
\item Prefer \textbf{fine-tuning} pre-trained models rather than training from scratch.
\item Use efficiency techniques: quantization, distillation, LoRA.
\item Report the carbon footprint in publications.
\item Choose data centers powered by renewable energy.
\end{itemize}
\end{bestpractice}

% ------------------------------------------------------------
\section{Regulatory Framework}
\label{sec:regulation}
\index{regulation}
\index{EU AI Act}

The \textbf{EU AI Act} (2024) classifies AI systems by risk level:
\begin{enumerate}
\item \textbf{Unacceptable risk}: prohibited (social scoring, subliminal manipulation).
\item \textbf{High risk}: regulated (healthcare, justice, recruitment).
\item \textbf{Limited risk}: transparency obligations.
\item \textbf{Minimal risk}: free use.
\end{enumerate}

% ------------------------------------------------------------
\section{Multimodal Models and Perspectives}
\label{sec:multimodal}
\index{multimodal}

\begin{sota}
\begin{itemize}
\item \textbf{Multimodal models}: GPT-4V, Claude 3.5 Vision, Gemini combine text, image, audio, and video in a single model.
\item \textbf{AI agents}: autonomous systems using tools (web browsing, code execution, API calls).
\item \textbf{World models}: learning physical representations for robotics and simulation (Sora, World Labs).
\item \textbf{Scientific AI}: accelerating discovery across all domains (materials, climate, biology).
\end{itemize}
\end{sota}

% ------------------------------------------------------------
\section{Chapter Summary}

\begin{keyformulas}
\begin{itemize}
\item \textbf{Detection}: YOLO predicts bounding boxes and classes per grid cell
\item \textbf{Segmentation}: U-Net with encoder-decoder and skip connections
\item \textbf{Grad-CAM}: $L^c = \relu(\sum_k \alpha_k^c A^k)$ with $\alpha_k^c = \frac{1}{Z}\sum_{ij}\frac{\partial y^c}{\partial A^k_{ij}}$
\item \textbf{SHAP}: Shapley values for attribution
\item \textbf{DP-SGD}: SGD with gradient clipping and Gaussian noise
\item \textbf{Federated learning}: weighted aggregation of local updates
\end{itemize}
\end{keyformulas}

% ------------------------------------------------------------
\section{Exercises}

\begin{exercise}[$\star$ --- Analysis]
Show that demographic parity and equalized odds cannot be satisfied simultaneously except in degenerate cases. What are the practical implications of this result?
\end{exercise}

\begin{exercise}[$\star$ --- Computation]
Compute the number of FLOPS required for a forward pass of a Transformer with $L=32$ layers, $d=4096$, $h=32$ heads, and a sequence length of $n=2048$. Estimate the energy cost on an A100 GPU.
\end{exercise}

\begin{exercise}[$\star\star$ --- Implementation]
Implement Grad-CAM for a ResNet-18 pre-trained on ImageNet. Visualize the attention maps for different classes and analyze the results.
\end{exercise}

\begin{exercise}[$\star\star$ --- Implementation]
Implement a simple federated learning pipeline with 5 clients on MNIST. Compare the performance with centralized training and analyze the impact of the number of communication rounds.
\end{exercise}

\begin{exercise}[$\star\star\star$ --- Research Project]
Conduct a fairness audit on a classification model (for example, income prediction on the Adult dataset). Compute demographic parity, equalized odds, and predictive equality. Propose and implement a debiasing method (pre-processing, in-training, or post-processing) and evaluate its impact on accuracy and fairness.
\end{exercise}
